\section{The \texttt{Monochromator\_pol} McStas Component}
Flat polarizaing monochromator crystal.

\subsection*{Identification}
\begin{itemize}
  \item \textbf{Site:} 
  \item \textbf{Author:} Peter Christiansen
  \item \textbf{Origin:} RISOE
  \item \textbf{Date:} 2006
\end{itemize}

\subsection*{Description}
\begin{lstlisting}
Based on Monochromator_flat.
Flat, infinitely thin mosaic crystal, useful as a monochromator or analyzer.
For an unrotated monochromator component, the crystal surface lies in the Y-Z
plane (ie. parallel to the beam).
The mosaic and d-spread distributions are both Gaussian.
Neutrons are just reflected (billard ball like). No correction is done for
mosaicity of reflecting crystal.
The crystal is assumed to be a ferromagnet with spin pointing up
eta-tilde = (0, 1, 0) (along y-axis), so that the magnetic field is
pointing opposite (0, -|B|, 0).

The polarisation is done by defining the reflectivity for spin up
(Rup) and spin down (Rdown) (which can be negative, see now!) and
based on this the nuclear and magnetic structure factors are
calculated:
FM = sign(Rup)*sqrt(|Rup|) + sign(Rdown)*sqrt(|Rdown|)
FN = sign(Rup)*sqrt(|Rup|) - sign(Rdown)*sqrt(|Rdown|)
and the physics is calculated as
Pol in = (sx_in, sy_in, sz_in)
Reflectivity R0 = FN*FN + 2*FN*FM*sy_in + FM*FM
(= |Rup| + |Rdown| (for sy_in=0))
Pol out:
sx = (FN*FN - FM*FM)*sx_in/R0;
sy = ((FN*FN - FM*FM)*sy_in + 2*FN*FM + FM*FM*sy_in)/R0;
sz = (FN*FN - FM*FM)*sz_in/R0;

These equations are taken from:
Lovesey: "Theory of neutron scattering from condensed matter, Volume
2", Eq. 10.96 and Eq. 10.110

This component works with gravity (uses PROP_X0).

Example: Monochromator_pol(zwidth=0.2, yheight=0.2, mosaic=30.0, dspread=0.0025, Rup=1.0, Rdown=0.0, Q=1.8734)

Monochromator lattice parameter
PG       002 DM=3.355 AA (Highly Oriented Pyrolythic Graphite)
PG       004 DM=1.677 AA
Heusler  111 DM=3.362 AA (Cu2MnAl)
CoFe         DM=1.771 AA (Co0.92Fe0.08)
Ge       111 DM=3.266 AA
Ge       311 DM=1.714 AA
Ge       511 DM=1.089 AA
Ge       533 DM=0.863 AA
Si       111 DM=3.135 AA
Cu       111 DM=2.087 AA
Cu       002 DM=1.807 AA
Cu       220 DM=1.278 AA
Cu       111 DM=2.095 AA
\end{lstlisting}

\subsection*{Input parameters}
Parameters in \textbf{boldface} are required; the others are optional.

\begin{longtable}{p{0.22\textwidth}p{0.12\textwidth}p{0.46\textwidth}p{0.14\textwidth}}
\toprule
\textbf{Name} & \textbf{Unit} & \textbf{Description} & \textbf{Default} \\
\midrule
\endhead
zwidth & m & Width of crystal & 0.1 \\
yheight & m & Height of crystal & 0.1 \\
mosaic &  & Mosaicity (FWHM) [arc minutes] & 30.0 \\
dspread & 1 & Relative d-spread (FWHM) & 0 \\
Q & AA-1 & Magnitude of scattering vector & 1.8734 \\
DM & AA & monochromator d-spacing instead of Q = 2*pi/DM & 0 \\
pThreshold &  & if probability\textgreater{}pThreshold then accept and weight else random & 0 \\
Rup & 1 & Reflectivity of neutrons with polarization up & 1 \\
Rdown & 1 & Reflectivity of neutrons with polarization down & 1 \\
debug &  & if debug \textgreater{} 0, print out some info about the calculations & 0 \\
\bottomrule
\end{longtable}

\subsection*{Links}
\begin{itemize}
  \item \href{run:/home/willend/willend-McCode/mcstas-comps/optics/Monochromator_pol.comp}{Source code} for \texttt{Monochromator\_pol.comp}.
\end{itemize}
\IfFileExists{Monochromator_pol_static.tex}{\input{Monochromator_pol_static.tex}}{}